\section{Solution Architecture}
\label{sec:architecture}

\subsection{The invariant beneath every surveyed platform}

Strip the branding from Tandem, TwinMaker, iTwin.js and Omniverse and the same three
layers appear. \textbf{Anchor}: a stable identifier tying a signal to a location on a
substrate --- IFC GUID, glTF node, USD prim, Element~ID, or a point or polygon in image
coordinates. \textbf{Encode}: a rule turning a value into a visual property --- colour,
opacity, visibility, transform, flow speed, ghost offset, glyph, label; this is where
commercial platforms put their product, and TwinMaker's no-code rule editor \emph{is}
what customers pay for. \textbf{Sustain}: applying encodings continuously at scale
without rebuilding the scene; iTwin.js devotes a native tile generator and a WebGL
frontend to this layer alone.

\begin{keypoint}
The renderer is not in this list. It is an implementation detail of the \emph{sustain}
layer. Treating it as the centre of the architecture is what makes a visualisation
package domain-locked --- and what makes packages unable to treat a photograph or a
video frame as a place where data can live.
\end{keypoint}

\subsection{Six layers}

Layers~1--4 are domain- and substrate-independent and form the reusable core; layer~5 is
a thin per-substrate adapter; layer~6 is third-party.

\paragraph{Layer 1 --- Transport.} Every source normalises to one event shape, with
optional historical backfill:

\begin{lstlisting}
export type Channel = 'measured' | 'simulated' | 'predicted' | 'setpoint';

export interface SignalSample {
  readonly twinId: string; readonly signalPath: string;   // "ward3/ahu1/supplyTemp"
  readonly channel: Channel; readonly ts: number;         // epoch ms, source-assigned
  readonly value: number | string | boolean | number[];
  readonly quality?: 'good' | 'stale' | 'bad';
}

export interface TransportAdapter {
  readonly id: string;
  connect(cfg: unknown): Promise<void>;
  subscribe(paths: string[], sink: (s: SignalSample) => void): () => void;
  range?(paths: string[], from: number, to: number): Promise<SignalSample[]>;
  disconnect(): Promise<void>;
}
\end{lstlisting}

Adapters follow the installed stack, not a wish-list: MQTT over WSS and AMQP over
Web-STOMP, both against the single RabbitMQ; ThingsBoard telemetry over its WebSocket
API; InfluxDB for \texttt{range()}; SQL/HTTP for the logger and arbitrary DT endpoints;
and a file adapter replaying CSV or Parquet from the library, GitLab or a workspace
URL --- which makes the stack testable offline and lets an FMI output file serve as the
\texttt{simulated} channel with no runtime integration.

\paragraph{Layer 2 --- Temporal state store, hot and cold.} A ring buffer per
\texttt{(signalPath, channel)} holds the live window; when the playhead moves outside
it, the store issues an InfluxDB \texttt{range()} query and serves from the returned
window. Every consumer calls \texttt{valueAt(path, channel, t)}, never ``current
value''. This implements \req{4} and unlocks 4D replay (\tech{6}), residuals (\tech{9}),
image-overlay replay (\tech{13}) and video sync (\tech{14}) at low marginal cost. The
playhead, subscriptions and connection status belong in Redux Toolkit; the sample
buffers deliberately do not, because dispatching an action per sample destroys frame
rate.

\begin{keypoint}
Having InfluxDB deployed converts the time store from a hard problem into a cache.
Without a historian, honest replay means holding everything in memory or building a
backend; with one, it is a windowed query behind an interface the rest of the package
cannot distinguish from live data.
\end{keypoint}

\paragraph{Layer 3 --- Anchor resolution.} Anchor kinds span three media, plus the
imported registry:

\begin{lstlisting}
export type AnchorKind =
  | 'ifc-guid' | 'gltf-node' | 'usd-prim' | 'instance-id'   // 3D substrates
  | 'geo'                                                    // georeferenced
  | 'image-point' | 'image-region' | 'image-plane'           // image substrates
  | 'video-panel'                                            // video substrates
  | 'tb-entity';                                             // ThingsBoard device/asset
\end{lstlisting}

An anchor carries \texttt{signalPath}, \texttt{kind}, \texttt{ref}, a human
\texttt{label} and, for image substrates, the \texttt{homographyId} of the calibration
that applies. \texttt{tb-entity} matters: because ThingsBoard holds the registry, an
anchor can pair \emph{this device} with \emph{that IFC element}, so the signal namespace
is imported rather than invented and a rename does not silently break the binding.

\paragraph{Layer 4 --- Encoding rules.} Declarative, serialisable, and the genuinely
reusable content of the package: \texttt{colorScale}, \texttt{visibility},
\texttt{transform}, \texttt{flow}, \texttt{residual}, \texttt{ghost}, \texttt{label} and
\texttt{attention} for geometric substrates; \texttt{glyph}, \texttt{regionFill},
\texttt{fieldOverlay} and \texttt{trajectory} for image and video substrates. Because
encodings are plain data, a complete visualisation is a JSON document --- which is what
makes \req{7} achievable and, given \pkg{@gitbeaker/rest}, what makes it storable and
reviewable with no new backend.

\paragraph{Layer 5 --- Substrate adapters.} The seam that delivers \req{1}. Generalising
from \emph{renderer} to \emph{substrate} is what admits images and video as peers of 3D:

\begin{lstlisting}
export interface SubstrateAdapter {
  readonly id: string; readonly supports: AnchorKind[];
  mount(container: HTMLElement, descriptor: SubstrateDescriptor): Promise<void>;
  resolve(anchor: Anchor): ElementRef | null;
  apply(ref: ElementRef, enc: ResolvedEncoding): void;   // hot path, must be O(1)
  frame(ref: ElementRef, opts?: FrameOptions): void;
  pick(x: number, y: number): ElementRef | null;
  setTime(t: number): void;                              // media seek, animation
  dispose(): void;
}
\end{lstlisting}

\begin{longtable}{@{}L{1.9cm}L{5.2cm}L{7.4cm}@{}}
\toprule
\textbf{Adapter} & \textbf{Backing library} & \textbf{Substrate and domains} \\
\midrule
\endhead
\texttt{image} & \pkg{react-konva}; \pkg{OpenCV.js} for calibration; \pkg{Leaflet}
(\texttt{CRS.Simple}) for very large images & Photographs, floor plans, scanned
drawings, rendered output frames. All domains; cheapest substrate, build first. \\
\addlinespace
\texttt{aec} & \pkg{@thatopen/components} + \pkg{@thatopen/fragments} & IFC models:
hospitals, buildings, plants. \\
\addlinespace
\texttt{mesh} & \pkg{react-three-fiber} + \pkg{drei} (+ \pkg{urdf-loader}) & glTF
assemblies: machines, robots, equipment. \\
\addlinespace
\texttt{video} & \texttt{HTMLVideoElement} + \pkg{hls.js} or WebRTC & Live and recorded
camera feeds beside or beneath simulation output. \\
\addlinespace
\texttt{geo} & \pkg{CesiumJS} / \pkg{3DTilesRendererJS} (+ \pkg{deck.gl}) & City,
mobility, marine; anything georeferenced or tiled. \\
\addlinespace
\texttt{field} & \pkg{vtk.js} & Simulation output: volumes, isosurfaces, streamlines.
Own canvas, synchronised camera. \\
\addlinespace
\texttt{embed} & \pkg{react-iframe} & Grafana panels, ThingsBoard dashboards and the
workspace VNC desktop, driven by the shared playhead. \\
\bottomrule
\end{longtable}

\subsection{Mapping onto the deployed DTaaS stack}
\label{sec:stack}

\begin{longtable}{@{}L{2.5cm}L{3.9cm}L{8.0cm}@{}}
\toprule
\textbf{Service} & \textbf{Role} & \textbf{How the package uses it} \\
\midrule
\endhead
\textbf{RabbitMQ} & Live transport, AMQP + MQTT from one broker & Two Layer-1 adapters
against one deployment: \pkg{mqtt.js} over WSS, \pkg{@stomp/stompjs} over Web-STOMP. No
second broker. \\
\addlinespace
\textbf{InfluxDB} & Historian & Cold tier of the Layer-2 store via
\pkg{@influxdata/influxdb-client}; makes \tech{6} and \tech{9} possible over arbitrary
past windows. \\
\addlinespace
\textbf{ThingsBoard} & Device/asset registry; telemetry & Anchor namespace
(\texttt{tb-entity}); a second live transport; dashboards embeddable. \\
\addlinespace
\textbf{Grafana} & Charting & \emph{Not} reimplemented. Embedded and driven by the
shared playhead: the viewer writes the current window into the panel URL's
\texttt{from}/\texttt{to}, so scrubbing the substrate moves the chart. \tech{10} for
near-zero cost. \\
\addlinespace
\textbf{GitLab} & Versioned storage and identity & Visualisation assets, anchor maps and
(via LFS) models and photographs read and written with \pkg{@gitbeaker/rest} under the
existing OAuth2 session. A threshold change becomes a merge request. \\
\addlinespace
\textbf{MongoDB} & Document store & Optional: large generated anchor maps, per-user view
state that should not pollute Git history. \\
\addlinespace
\textbf{Library / runner} & Files; DT execution & Resolves \texttt{lib://} substrates
(needs MIME types for \texttt{.glb}, \texttt{.ifc}, \texttt{.frag}); produces the
\texttt{simulated} channel and rendered frames. \\
\bottomrule
\end{longtable}

\begin{keypoint}
The highest-value integration in this table is the cheapest: driving an embedded Grafana
panel from the visualisation playhead. It gives linked time-and-space navigation without
writing a charting component, and respects the division of labour the marine DT pilot
arrived at independently~\citep{vasilijevic2024trondheim}.
\end{keypoint}

\subsection{Persistence, the workspace, and composability}
\label{sec:workspace}

Because \pkg{@gitbeaker/rest} is already a client dependency and GitLab is already the
identity provider, the visualisation asset needs no bespoke storage: the viewer reads
\texttt{visualisation.json} and its substrate at a given ref, the binding editor commits
changes back, review happens as a merge request diff, and large substrates live in Git
LFS beside the configuration referencing them.

Some visualisation will never happen in a browser canvas --- full CFD post-processing,
ParaView, Blender, \texttt{rviz}, a Unity city build. All already run in the DTaaS
workspace, whose xfce desktop is reachable over VNC inside the client. The architecture
should embrace this on three conditions: the desktop mounts through the \texttt{embed}
adapter like any other substrate; output files written to the workspace come back by URL
and are consumed by the \texttt{image} or \texttt{mesh} adapter, so a heavyweight render
becomes a navigable, time-aligned artefact for everyone else; and where a workspace tool
cannot accept the client's playhead, the layout says the pane is independent rather than
implying synchrony. This gives an honest answer to ``can DTaaS do proper scientific
visualisation?'' --- yes, in the workspace --- without pretending the web layer will
match ParaView.

Finally, DTaaS already treats data, models, functions and tools as composable library
assets. The consistent move is to add a \emph{visualisation} asset type, so a
\texttt{visualisation.json} sits beside the models and functions of a twin and is picked
in the same composition flow. An abbreviated example spanning three substrates:

\begin{lstlisting}
{ "schemaVersion": "1.1", "name": "junction-7-live",
  "layout": { "type": "split", "panes": ["photo", "camera", "grafana"] },
  "substrates": {
    "photo":   { "adapter": "image", "source": "git://twins/junction7/junction7.jpg",
                 "calibration": { "type": "homography", "id": "h1",
                   "imagePoints": [[412,880],[1502,872],[1710,1040],[210,1050]],
                   "worldPoints": [[0,0],[24,0],[24,12],[0,12]], "units": "m" } },
    "camera":  { "adapter": "video", "source": ".../junction7/index.m3u8",
                 "latencyHintMs": 2500 },
    "grafana": { "adapter": "embed", "source": "https://grafana.example/d/abc/junction7",
                 "timeParams": ["from", "to"] } },
  "transports": [
    { "adapter": "mqtt", "url": "wss://rabbitmq.example:15676/ws",
      "topics": ["dtaas/junction7/+/count"], "channel": "measured" },
    { "adapter": "influx", "org": "dtaas", "bucket": "junction7", "role": "history" },
    { "adapter": "file", "url": "lib://data/junction7-sumo-out.csv",
      "channel": "simulated" } ],
  "anchors": [
    { "signalPath": "junction7/armA/queueLen", "kind": "image-region",
      "ref": "poly:armA", "homographyId": "h1", "label": "North arm queue" } ],
  "encodings": [
    { "target": "junction7/armA/queueLen", "substrate": "photo",
      "encoding": { "type": "regionFill", "domain": [0,40], "scheme": "traffic" } },
    { "target": "junction7/armA/queueLen", "substrate": "photo",
      "encoding": { "type": "residual", "against": "simulated", "domain": [-8,8] } } ] }
\end{lstlisting}

A visualisation therefore becomes shareable (replace the photograph and the homography),
reviewable (JSON in Git) and executable by the runner (a headless adapter emits the same
overlay as an image file during a DT execution).

\subsection{Design rules}

\begin{enumerate}[itemsep=1pt]
  \item The binding layer is the product; substrates are dependencies. Layers~1--4 must
        import no rendering, imaging or media library --- enforce with a lint rule.
  \item One clock. Scene, image overlay, video element and embedded Grafana panel all
        sample the playhead.
  \item Encodings are data. Anything expressible only as code is a gap in the
        vocabulary.
  \item Hot path outside Redux. Redux holds control state; per-frame sampling reads
        buffers directly.
  \item Reuse before build. If InfluxDB, RabbitMQ, ThingsBoard, Grafana or GitLab does
        it, call it.
  \item Every adapter lazily loaded. An image twin must never ship the BIM engine.
\end{enumerate}

